\documentclass[11pt]{amsart}

\usepackage{amsmath,amssymb,amsthm}
\usepackage[margin=1in]{geometry}
\usepackage{enumitem}

\newtheorem{theorem}{Theorem}[section]
\newtheorem{proposition}[theorem]{Proposition}
\newtheorem{lemma}[theorem]{Lemma}
\newtheorem{corollary}[theorem]{Corollary}
\theoremstyle{definition}
\newtheorem{definition}[theorem]{Definition}
\theoremstyle{remark}
\newtheorem{remark}[theorem]{Remark}

\newcommand{\M}{\mathcal M}
\newcommand{\overlineM}{\overline{\mathcal M}}
\newcommand{\C}{\mathbb C}
\newcommand{\DD}{\Delta}
\newcommand{\cV}{\mathcal V}
\newcommand{\cC}{\mathcal C}
\newcommand{\norm}{\mathcal N}
\newcommand{\defo}{\mathrm{Def}}

\begin{document}

\title{Non-variation of fixed boundary factors\\
for deformations of a totally geodesic surface}
\date{}
\maketitle

\section{Statement}

Let \(V_{0}\subset \M_g\) be a smooth complex surface which is totally
geodesic for the Teichm\"uller metric, and let
\[
\cV\longrightarrow \DD,\qquad V_t=\cV_t,
\]
be a holomorphic deformation of \(V_0\) satisfying the smooth-triviality
and boundary-compatibility assumptions (i) and (ii) in the question.

Fix a boundary stratum \(\Delta_\Gamma\) which, after passing to a finite
\^etale chart, is a product
\[
\Delta_\Gamma
   =\prod_{j=1}^{k}\M_{g_j,n_j},
\]
and let
\[
\pi_i:\Delta_\Gamma\longrightarrow \M_{g_i,n_i}
\]
be projection to one of the factors. Put
\[
\partial V_t=\overline{V_t}\cap \Delta_\Gamma .
\]

We prove the following.

\begin{theorem}\label{thm:main}
Assume that, for every \(t\in\DD\),
\begin{enumerate}[label=\textup{(\alph*)}]
\item \(\dim_\C\partial V_t=1\);
\item there is a point \(Y_t\in\M_{g_i,n_i}\) such that
\[
\pi_i(\partial V_t)=\{Y_t\}.
\]
\end{enumerate}
Then
\[
Y_t=Y_0\qquad\text{for every }t\in\DD.
\]
Equivalently,
\[
\pi_i(\partial V_t)=\{Y_0\}
\qquad\text{for every }t.
\]
\end{theorem}

The other factors of \(\partial V_t\) may move. Only the factor on which
the projection is a singleton is rigid.

\section{Geometric preparation}

All local arguments may be made on finite \'etale charts, so we suppress
the harmless orbifold quotients. This is permitted by the assumption in
the question and does not affect constancy of the projection.

\subsection{The total boundary family}

Let
\[
\cC_\Gamma\subset \overline{\cV}
\]
be the closure of the locus lying over \(\Delta_\Gamma\). Thus
\[
(\cC_\Gamma)_t=\partial V_t .
\]
By assumption (i), the morphism
\[
\cC_\Gamma\longrightarrow\DD
\]
is a smoothly trivial family of smooth complex curves. In particular, for
every \(x_0\in\partial V_0\), there are a neighborhood
\(W\subset\cC_\Gamma\) of \((x_0,0)\) and a holomorphic section
\[
t\longmapsto x(t)\in\partial V_t,\qquad x(0)=x_0.
\]
Set
\[
y(t):=\pi_i(x(t)).
\]
Since \(\pi_i(\partial V_t)\) is the singleton \(\{Y_t\}\), this is
independent of the chosen section and
\[
y(t)=Y_t.
\]
Hence \(t\mapsto Y_t\) is holomorphic.

\subsection{Totally geodesic central fibres and affine period models}

We use the following standard structure theorem.

\begin{theorem}[M\"oller]\label{thm:Moller}
Every complex totally geodesic subvariety of \(\M_g\) is, after lifting to
a suitable stratum of the Hodge bundle and passing to a finite cover, the
projection of an affine invariant submanifold. In local period
coordinates, its lift is cut out by complex-linear equations.
\end{theorem}

This is the main theorem of M\"oller,
\emph{Teichm\"uller curves and totally geodesic subvarieties}
\cite[Theorem~1.1 and \S4]{Moller2008}. The underlying
orbit-closure theory is due to Eskin--Mirzakhani--Mohammadi
\cite{EMM2015}.

The boundary of such an affine invariant manifold is also affine in the
sense made precise by Mirzakhani and Wright.

\begin{theorem}[Mirzakhani--Wright]\label{thm:MWboundary}
The boundary of an affine invariant submanifold in a product of lower
strata is, up to the usual finite covers, itself an affine invariant
submanifold. Moreover, incidence-preserving boundary deformations are
identified with the boundary affine tangent space.
\end{theorem}

This is \cite[Theorem~1.1 and the boundary-tangent-space construction in
\S\S2--4]{MW2017}. Cylinder deformations give the integration of the
relevant boundary tangent vectors; see Wright
\cite[Theorem~1.1]{Wright2015}.

\subsection{The boundary normal cone}

Let \(C_0\subset\partial V_0\) be a local component and let \(N_0\) be a
lift of \(V_0\) to an affine period chart, so that the lift of \(C_0\) is
a boundary component \(B_0\) of \(N_0\). At \(x\in B_0\), write
\[
\operatorname{pr}_i:
T_x\Delta_\Gamma\longrightarrow T_{\pi_i(x)}\M_{g_i,n_i}
\]
for the differential of the projection to the \(i\)-th factor.

Since \(\pi_i(C_0)=\{Y_0\}\), one has
\begin{equation}\label{eq:zero-boundary-projection}
\operatorname{pr}_i(T_xB_0)=0
\qquad (x\in B_0).
\end{equation}

We shall use the following boundary rigidity criterion. It is a standard
consequence of Theorems~\ref{thm:Moller} and~\ref{thm:MWboundary}; we
isolate it because it is the only place where total geodesicity is used.

\begin{lemma}[Boundary normal rigidity]
\label{lem:boundary-normal-rigidity}
Let \(\sigma\) be an infinitesimal normal deformation of \(V_0\) which is
compatible with every boundary stratum in the sense of assumptions
\textup{(i)} and \textup{(ii)}. Suppose that \(\sigma\) is tangent to
\(\Delta_\Gamma\) along \(C_0\). Then
\[
\operatorname{pr}_i(\sigma|_{C_0})=0.
\]
\end{lemma}

\begin{proof}
Pass to a finite cover on which Theorem~\ref{thm:Moller} applies, and let
\(N_0\) and \(B_0\) be as above.

The period-coordinate description in Theorem~\ref{thm:Moller} gives a
canonical flat reduction of the logarithmic normal sheaf of \(N_0\). A
normal deformation compatible with all boundary strata is horizontal for
this reduction; this is the infinitesimal form of M\"oller's affine
linearization theorem \cite[\S4]{Moller2008}. Consequently
\(\operatorname{pr}_i(\sigma|_{B_0})\) is a parallel vector field along
the affine boundary component \(B_0\).

Suppose that \(\operatorname{pr}_i(\sigma|_{B_0})\neq0\). By the
Mirzakhani--Wright boundary theorem, this parallel factor direction is a
boundary deformation direction of the affine invariant boundary component.
Wright's cylinder-deformation theorem integrates such a boundary
direction: an arc of points of \(B_0\) remains in the boundary of the same
affine invariant manifold while its projection to the \(i\)-th factor
moves in the direction of this vector field
\cite{MW2017,Wright2015}. It follows that
\(\operatorname{pr}_i(B_0)\) contains an arc based at \(Y_0\). This
contradicts
\[
\operatorname{pr}_i(B_0)=\{Y_0\}.
\]
Therefore \(\operatorname{pr}_i(\sigma|_{B_0})=0\). Descending the finite
cover gives the lemma.
\end{proof}

\begin{remark}
Assumptions (i) and (ii) are essential here. They ensure that a
deformation of \(V_0\) induces a logarithmic normal section, rather than a
section with uncontrolled poles at the boundary, and that deformation in a
larger stratum is compatible with its closures in smaller strata.
\end{remark}

\section{Leading jets of a moving boundary point}

Choose a local holomorphic chart on \(\M_{g_i,n_i}\) around \(Y_0\), and
continue to write
\[
Y_t=(y_1(t),\ldots,y_m(t)),\qquad Y_0=(0,\ldots,0).
\]
Suppose that \(Y_t\) is not identically equal to \(Y_0\). Then there are
an integer \(r\geq1\) and a nonzero vector \(w\in T_{Y_0}\M_{g_i,n_i}\)
such that
\begin{equation}\label{eq:leading-term}
Y_t=Y_0+t^r w+O(t^{r+1}).
\end{equation}
Here \(w\) is interpreted through the chosen local chart.

\begin{lemma}[Leading-jet lemma]\label{lem:leading-jet}
Under the hypotheses of Theorem~\ref{thm:main}, if \eqref{eq:leading-term}
holds, then there is a compatible infinitesimal normal deformation
\(\sigma\) of \(V_0\) such that
\[
\operatorname{pr}_i(\sigma|_{C_0})=w
\]
as a constant vector field along \(C_0\).
\end{lemma}

\begin{proof}
This is the standard leading-coefficient construction in embedded
deformation theory. Locally near a point of \(C_0\), choose holomorphic
coordinates on \(\overlineM_g\times\DD\) adapted to the boundary divisor
and write the defining equations of the total family as
\[
F_a(z,t)=0.
\]
Expand
\[
F_a(z,t)=\sum_{q\geq0}t^qF_{a,q}(z).
\]
The first order at which the \(i\)-th boundary factor moves is \(r\).
Differentiating the incidence equations \(F_a(z,t)=0\) \(r\) times, all
terms involving the \(i\)-th coordinate of orders \(<r\) vanish by the
choice of \(r\). The remaining coefficient defines a Cech \(0\)-cocycle
with values in the normal sheaf of the central fibre.

Assumption (i) says that these local cocycles glue over the smooth trivial
family \(C_0\). Assumption (ii) says that the resulting section has no
poles along smaller boundary strata, hence lies in the compatible
logarithmic normal sheaf. Its projection to the tangent space of the
\(i\)-th factor is precisely the leading coefficient \(w\).

This is the analytic version of the usual leading-term construction for
embedded deformations; see, for example, \cite[Chapter~1]{Sernesi2006}.
\end{proof}

\section{Proof of Theorem~\ref{thm:main}}

Assume, for contradiction, that \(t\mapsto Y_t\) is not constant. By
\eqref{eq:leading-term}, after choosing local coordinates around \(Y_0\),
there exist \(r\geq1\) and \(w\neq0\) such that
\[
Y_t=Y_0+t^r w+O(t^{r+1}).
\]

By Lemma~\ref{lem:leading-jet}, there is a compatible infinitesimal
normal deformation \(\sigma\) of \(V_0\) whose restriction to
\(C_0\subset\partial V_0\) projects to the constant vector \(w\) on
\(\M_{g_i,n_i}\):
\[
\operatorname{pr}_i(\sigma|_{C_0})=w.
\]

On the other hand, Lemma~\ref{lem:boundary-normal-rigidity}, applied to
the totally geodesic central fibre \(V_0\), gives
\[
\operatorname{pr}_i(\sigma|_{C_0})=0.
\]
Thus \(w=0\), a contradiction.

Consequently every coefficient in the Taylor expansion of \(Y_t-Y_0\) at
\(t=0\) vanishes. Since \(Y_t\) depends holomorphically on \(t\), the
identity theorem yields
\[
Y_t=Y_0
\]
for all \(t\) in the connected unit disk. Hence
\[
\boxed{\ \pi_i(\partial V_t)=\{Y_0\}\quad\text{for all }t\in\DD.\ }
\]

This completes the proof.

\begin{thebibliography}{99}

\bibitem[BCGGM2018]{BCGGM2018}
M.~Bainbridge, D.~Chen, Q.~Gendron, S.~Grushevsky, and M.~M\"oller,
\newblock Compactification of strata of Abelian differentials,
\newblock \emph{Duke Math. J.} \textbf{167} (2018), no.~12, 2347--2416.

\bibitem[EMM2015]{EMM2015}
A.~Eskin, M.~Mirzakhani, and A.~Mohammadi,
\newblock Isolation, equidistribution, and orbit closures for the
\(\mathrm{SL}(2,\mathbb R)\) action on moduli space,
\newblock \emph{Ann. of Math. (2)} \textbf{182} (2015), no.~2, 673--721.

\bibitem[Filip2016]{Filip2016}
S.~Filip,
\newblock Splitting mixed Hodge structures over affine invariant manifolds,
\newblock \emph{Ann. of Math. (2)} \textbf{183} (2016), no.~2, 681--713.

\bibitem[M\"oller2008]{Moller2008}
M.~M\"oller,
\newblock Teichm\"uller curves and totally geodesic subvarieties,
\newblock \emph{Geom. Funct. Anal.} \textbf{18} (2008), no.~4, 1135--1169.

\bibitem[Mirzakhani--Wright2017]{MW2017}
M.~Mirzakhani and A.~Wright,
\newblock The boundary of an affine invariant submanifold,
\newblock \emph{Invent. Math.} \textbf{209} (2017), no.~3, 927--984.

\bibitem[Sernesi2006]{Sernesi2006}
E.~Sernesi,
\newblock \emph{Deformations of Algebraic Schemes},
\newblock Grundlehren der mathematischen Wissenschaften \textbf{334},
Springer-Verlag, Berlin, 2006.

\bibitem[Wright2015]{Wright2015}
A.~Wright,
\newblock Cylinder deformations in orbit closures of translation surfaces,
\newblock \emph{Geom. Topol.} \textbf{19} (2015), no.~1, 413--438.

\end{thebibliography}

\end{document}